\documentclass[10pt, letterpaper]{article}

% Packages:
\usepackage[
    ignoreheadfoot, % set margins without considering header and footer
    top=0.1 cm, % seperation between body and page edge from the top
    bottom=0.2 cm, % seperation between body and page edge from the bottom
    left=2 cm, % seperation between body and page edge from the left
    right=2 cm, % seperation between body and page edge from the right
    footskip=1.0 cm, % seperation between body and footer
    % showframe % for debugging 
]{geometry} % for adjusting page geometry
\usepackage{titlesec} % for customizing section titles
\usepackage{tabularx} % for making tables with fixed width columns
\usepackage{array} % tabularx requires this
\usepackage[dvipsnames]{xcolor} % for coloring text
\definecolor{primaryColor}{RGB}{0, 0, 0} % define primary color
\usepackage{enumitem} % for customizing lists
\usepackage{fontawesome5} % for using icons
\usepackage{amsmath} % for math
\usepackage[
    pdftitle={Aiden Wang's 10/18/2025 CV},
    pdfauthor={Aiden Wang},
    pdfcreator={LaTeX with RenderCV},
    colorlinks=true,
    urlcolor=primaryColor
]{hyperref} % for links, metadata and bookmarks
\usepackage[pscoord]{eso-pic} % for floating text on the page
\usepackage{calc} % for calculating lengths
\usepackage{bookmark} % for bookmarks
\usepackage{lastpage} % for getting the total number of pages
\usepackage{changepage} % for one column entries (adjustwidth environment)
\usepackage{paracol} % for two and three column entries
\usepackage{ifthen} % for conditional statements
\usepackage{needspace} % for avoiding page brake right after the section title
\usepackage{iftex} % check if engine is pdflatex, xetex or luatex

% Ensure that generate pdf is machine readable/ATS parsable:
\ifPDFTeX
    \input{glyphtounicode}
    \pdfgentounicode=1
    \usepackage[T1]{fontenc}
    \usepackage[utf8]{inputenc}
    \usepackage{lmodern}
\fi

\usepackage{charter}

% Some settings:
\raggedright
\AtBeginEnvironment{adjustwidth}{\partopsep0pt} % remove space before adjustwidth environment
\pagestyle{empty} % no header or footer
\setcounter{secnumdepth}{0} % no section numbering
\setlength{\parindent}{0pt} % no indentation
\setlength{\topskip}{0pt} % no top skip
\setlength{\columnsep}{0.15cm} % set column seperation
\pagenumbering{gobble} % no page numbering

\titleformat{\section}{\needspace{4\baselineskip}\bfseries\large}{}{0pt}{}[\vspace{1pt}\titlerule]

\titlespacing{\section}{
    % left space:
    -1pt
}{
    % top space:
    0.1 cm
}{
    % bottom space:
    0.1 cm
} % section title spacing

\renewcommand\labelitemi{$\vcenter{\hbox{\small$\bullet$}}$} % custom bullet points
\newenvironment{highlights}{
    \begin{itemize}[
        topsep=0.10 cm,
        parsep=0.10 cm,
        partopsep=0pt,
        itemsep=0pt,
        leftmargin=0 cm + 10pt
    ]
}{
    \end{itemize}
} % new environment for highlights


\newenvironment{highlightsforbulletentries}{
    \begin{itemize}[
        topsep=0.10 cm,
        parsep=0.10 cm,
        partopsep=0pt,
        itemsep=0pt,
        leftmargin=10pt
    ]
}{
    \end{itemize}
} % new environment for highlights for bullet entries

\newenvironment{onecolentry}{
    \begin{adjustwidth}{
        0 cm + 0.00001 cm
    }{
        0 cm + 0.00001 cm
    }
}{
    \end{adjustwidth}
} % new environment for one column entries

\newenvironment{twocolentry}[2][]{
    \onecolentry
    \def\secondColumn{#2}
    \setcolumnwidth{\fill, 4.5 cm}
    \begin{paracol}{2}
}{
    \switchcolumn \raggedleft \secondColumn
    \end{paracol}
    \endonecolentry
} % new environment for two column entries

\newenvironment{threecolentry}[3][]{
    \onecolentry
    \def\thirdColumn{#3}
    \setcolumnwidth{, \fill, 4.5 cm}
    \begin{paracol}{3}
    {\raggedright #2} \switchcolumn
}{
    \switchcolumn \raggedleft \thirdColumn
    \end{paracol}
    \endonecolentry
} % new environment for three column entries

\newenvironment{header}{
    \setlength{\topsep}{0pt}\par\kern\topsep\centering\linespread{1.5}
}{
    \par\kern\topsep
} % new environment for the header

\newcommand{\placelastupdatedtext}{% \placetextbox{<horizontal pos>}{<vertical pos>}{<stuff>}
  \AddToShipoutPictureFG*{% Add <stuff> to current page foreground
    \put(
        \LenToUnit{\paperwidth-2 cm-0 cm+0.05cm},
        \LenToUnit{\paperheight-1.0 cm}
    ){\vtop{{\null}\makebox[0pt][c]{
        \small\color{gray}\textit{Last updated in September 2024}\hspace{\widthof{Last updated in September 2024}}
    }}}%
  }%
}%

% save the original href command in a new command:
\let\hrefWithoutArrow\href

% new command for external links:


\begin{document}

\begin{center}
    {\LARGE \textbf{Chenshuo (Aiden) Wang}}\\[3pt]
    \href{mailto:chens001@umn.edu}{chens001@umn.edu} \textbar\ 
    \href{https://linkedin.com/in/aiden-wang-964082354}{linkedin.com/in/aiden-wang} \textbar\ 
    \href{https://github.com/altoauto}{github.com/altoauto}
\end{center}
\vspace{-5mm}

\section*{Education}
\textbf{University of Minnesota – Twin Cities} — B.S. Mechanical Engineering \hfill May 2026 \,|\, GPA: 3.4/4.0 \\
\textbf{Northeastern University} — M.S. Energy Systems \hfill Expected May 2028

\section{Experience}
\textbf{AltoAuto} --- \textit{Founder / Engineering Automation Developer} \hfill 2024 -- Present\\
\vspace{0.1mm}
\textit{Developed CAD--FEA automation workflows for engineering teams.}
\begin{itemize}[leftmargin=1.5em, itemsep=2pt, topsep=2pt]
    \item Architected an end-to-end pipeline integrating \textbf{SolidWorks API}, \textbf{Gmsh}, \textbf{PyMAPDL}, and \textbf{Code\_Aster} to automate meshing, boundary condition mapping, and solver setup.
    \item Created a parametric batch framework for reproducible studies and validation, reducing model preparation time by 40\%.
    \item Automated stress visualization, convergence checks, and report-ready outputs across solver platforms using a PyVista-based post-processing pipeline.
    \item Contributed to the PyMAPDL open-source community by testing visualization modules and incorporating feedback into extensions.
\end{itemize}

\vspace{+1.5mm}

\textbf{Daikin Industries} —\textit{ Lead Modeling \& Controls Engineer (Sponsored Project)}\hfill 2025 -- 2026\\[-2pt]
\vspace{0.1mm}
\textit{Industry-sponsored embedded controls and system modeling for valve test bench.}
\begin{itemize}[leftmargin=1.5em, itemsep=2pt, topsep=2pt]
    \item Designed and tuned PI feedback control system; validated closed-loop stability, disturbance rejection, and transient response across defined operating envelopes and safety constraints.
    \item Implemented modular parameter-sweep and batch-run framework ($UA$, $\Delta P$, NTU) to accelerate calibration iteration and support validation against hardware test bench requirements.
    \item Collaborated directly with Daikin controls and MBD engineers to validate open-loop and closed-loop system behavior against test specifications and safety envelopes.
\end{itemize}
\vspace{+0.5mm}


\section{Technical Projects}
\textbf{Industrial Robotics \& Workcell Automation} — \textit{UR5 Manufacturing Workcell} \\[-2pt]
\vspace{0.1mm}
\begin{itemize}[leftmargin=1.5em,itemsep=2pt,topsep=2pt]
    \item Developed Python and RoboDK automation routines for pick-and-place operations, enforcing hardware-safe joint staging, kinematic limits, and saturation-aware command sequencing for repeatable production motion.
    \item Established complete workcell transform chain; applied FK/IK with consistent configuration branches to guarantee reachable, repeatable end-effector poses --- validated through structured commissioning tests directly analogous to FAT/SAT protocols.
\end{itemize}
\vspace{+0.5mm}

\vspace{+0.5mm}
\textbf{Motion Control Project} — \textit{DC Motor System ID \& Model-Based Closed-Loop Control} \\[-2pt]
\vspace{0.1mm}
\begin{itemize}[leftmargin=1.5em,itemsep=2pt,topsep=2pt]
    \item Engineered a 750\,Hz real-time data acquisition and actuation control loop in C/C++, integrating ADC/DAC components, quadrature encoders, and signal conditioning for hardware-in-the-loop testing.
    \item Executed dynamic excitation tests to derive a discrete-time plant model, characterizing and mitigating non-ideal electromechanical regimes including deadzone stiction, sensor quantization, and torque saturation.
    \item Deployed a model-based PID control architecture for precision speed and position tracking; automated performance validation and disturbance rejection analysis using a custom Python/MATLAB data pipeline.
\end{itemize}
\vspace{+0.5mm}

\textbf{CFD \& Multiphase Flow} — \textit{ACDI Interface Tracking \& Turbulence Modeling} \\[-2pt]
\vspace{0.1mm}
\begin{itemize}[leftmargin=1.5em,itemsep=2pt,topsep=2pt]
    \item Implemented CDI and ACDI phase-field methods (Jain 2022) in Python for multiphase interface tracking; benchmarked upwind, central, and skew-symmetric flux schemes with Euler and RK4 integration across drop advection and shear-flow test cases.
    \item Developed a 2D incompressible Navier--Stokes solver coupled with ACDI to simulate droplet impact ($Re=500$, $We=100$) using fractional-step projection and continuum surface force for surface tension modeling.
    \item Built a 2D Finite Volume $k$--$\omega$ SST solver with SIMPLE pressure--velocity coupling; validated velocity profiles and wall shear stress against channel-flow benchmarks.
\end{itemize}


\section*{Technical Skills}
\textbf{Programming \& Automation:} Python, C/C++, MATLAB/Simulink, Git, Arduino, LabVIEW, Allen-Bradley ControlLogix, Ladder Logic, PLC Programming, PID/PI Control  \\
\textbf{Modeling \& Simulation:} Simscape, ANSYS (Mechanical \& Fluent), RoboDK, Code\_Aster, Gmsh, OpenFOAM, OpenModelica, SolidWorks, Creo \\
\textbf{Domains:} Computational Design \& Simulation (CAE), Thermal Management Systems, Mechatronics \& Control, Design Process Automation 

\end{document}